\section{Unconditional quotient returns at the critical-strip endpoints}
\label{sec:endpoint-quotient-return}

The classical zero-free region supplies an arithmetic instance
of the polynomial-cost criterion at the exact endpoints of
the critical strip.  This instance includes every fixed
polynomial weight.  It does not concern a fixed narrower
strip whose edges lie strictly between zero and one.

Fix a primitive finite-order Hecke character over a fixed
number field, allowing the principal character, and use its
same pole-removed entire completion \(\Xi\), ideal,
adapted heights and holomorphic block norms as in
MT-M18 and (MA-PR.1).  For any fixed \(a\in\mathbb R\),
put
\[
 \begin{gathered}
 u=-a,\qquad v=1-a,\qquad s\ge0,\\
 H=L^2\!\left(\mathbb R,
       \max(e^{-2ax},e^{2(1-a)x})(1+x^2)^{s/2}\,dx\right).
 \end{gathered}
 \tag{MA-PR.200}
\]
Thus \(e_-=a+u=0\) and \(e_+=a+v=1\).
In particular \(a=1/2\) is the centered endpoint
\(\alpha=1/2\).  The quotient \(Q_H\), the map \(i\),
the finite restrictions \(\mathscr R_n\), and \(J_n\)
retain exactly their original definitions and norms.

\paragraph{MA-PR-014 (the endpoint polynomial bound and actual returns).}
For the fixed data in (MA-PR.200), every completed zero
is retained in full and there are finite constants
\(C,D\) such that
\[
       \mathcal C_n
        =\|\iota_n\mathscr R_n\|_{\mathcal L(Q_H,J_n)}
                    \le C(2+|n|)^D
                         \qquad(n\in\mathbb Z).
 \tag{MA-PR.201}
\]
This is the norm on the whole Hilbert quotient, with
the original full holomorphic quotient as target.
Consequently every \(g\in\mathcal E\) has the actual
unique equivariant quotient return
\[
 \begin{aligned}
 S_g:Q_H&\longrightarrow Q_\infty,\\
 iS_g&=C_g^{Q_H},\\
 S_gU_t^{Q_H}&=e^{-at}U_t^{Q_\infty}S_g,\\
 S_gi&=C_{T_{-a}g}^{Q_\infty}.
 \end{aligned}
 \tag{MA-PR.202}
\]
Its \(\mathcal E\)-valued lifts can be chosen jointly
continuous and bilinear in \((g,v)\), with the common
all-order nuclear representations of MA-PR-010,
uniformly when \(g\) varies in a bounded subset of
\(\mathcal E\).
In particular this applies to all compact smooth
tests, including every \(g\) with \(g(0)=0\) and the
entire space \(\mathcal D_{[1,2]}\).

The corresponding quotient traces satisfy
\[
       \operatorname{Tr}_{Q_H} C_g
        =\operatorname{Tr}_{Q_\infty}
                                  C_{T_{-a}g}.
 \tag{MA-PR.203}
\]
No A, M, GRH, simplicity, or uniform separation
between distinct zeros is assumed.  Constants are
for this fixed character, field and weight; no
uniformity over an arithmetic family is asserted.

\begin{proof}
\emph{1. The two established arithmetic estimates.}
The classical zero-free region and the functional
equation give a constant \(c>0\), for the fixed
function, such that every completed zero
\(\rho=\beta+i\gamma\) satisfies
\[
        \min(\beta,1-\beta)
                    \ge \frac{c}{\log(3+|\gamma|)}.
 \tag{MA-PR.204}
\]
Here the high-ordinate zero-free estimate is
Lagarias, Montgomery, and Odlyzko (1979),
Lemma~2.3, (2.16), with proof on pp.~277--279.
Apply it also to the dual character and use
\(\rho\mapsto1-\rho\) in the functional equation
to obtain the left-edge estimate.
The completed divisor lies strictly in \(0<\beta<1\).
Only finitely many zeros have bounded ordinate.
Their positive distances from the two edges,
including any exceptional real zero and its dual
reflection, can therefore be absorbed by decreasing
\(c\).  This is a fixed-function statement and
requires no uniform lower bound on exceptional
zeros in a varying family.

The radius-sensitive classical count is
\[
 \begin{split}
 N_e(t,r)
   &:=
      \sum_{\substack{\Xi(\rho)=0\\|\rho-e-it|<r}}
                                         m_\rho\\
   &\le C_0+C_1r\log(3+|t|),
       \qquad e\in\{0,1\},\quad 0<r\le1 .
 \end{split}
 \tag{MA-PR.205}
\]
At \(e=1\) this is Lagarias, Montgomery, and
Odlyzko (1979), Lemma~2.2, (2.11)--(2.14),
pp.~276--277.  Although its statement uses centers
with real part greater than one, its proof explicitly
establishes the estimate at \(1+it\).  Alternatively,
use centers \(1+\epsilon+it\), enlarge the radius
by \(\epsilon\), and let \(\epsilon\downarrow0\);
the bound has constants uniform in that displacement.
The left-edge count follows by applying the
right-edge count to the dual character at height
\(-t\).  Multiplicities are preserved in both
steps.  An explicit version of the same count is
Thorner and Zaman (2015), Lemma~2.5, printed p.~9.
These established arithmetic results are used as
cited; the argument below proves their consequence
for the actual quotient cost.

\emph{2. The complex-distance potential is logarithmic.}
Write \(\nu_n=2+|n|\) and
\(\ell_n=\log(e\nu_n)\).
For \(\rho\in\mathcal Z_n\) of (MA-PR.153),
\(3+|\gamma|\) and \(\nu_n\) are comparable.
Equation (MA-PR.204) therefore gives, after
decreasing a fixed \(c_0>0\),
\[
   |\rho-e-it|\ge|\beta-e|
                  \ge c_0/\ell_n
       \quad(e=0,1,\ \rho\in\mathcal Z_n).
 \tag{MA-PR.206}
\]
Choose \(c_0\) also so that \(c_0/\ell_n\le1/2\).
For \(h_n\le t\le h_{n+1}\), the elementary
finite-sum integral identity gives
\[
 \begin{split}
 &\sum_{\rho\in\mathcal Z_n}m_\rho
                  \log^+\frac1{|\rho-e-it|}\\
 &\qquad
 =\int_0^1
       \left(\sum_{\substack{\rho\in\mathcal Z_n\\
                               |\rho-e-it|<r}}m_\rho\right)
                                                      \frac{dr}{r}\\
 &\qquad
 \le \int_{c_0/\ell_n}^{1}
                         (C_0+C_2r\ell_n)\frac{dr}{r}\\
 &\qquad
 =C_0\log(\ell_n/c_0)+C_2(\ell_n-c_0)
                                      \le C_3\ell_n .
 \end{split}
 \tag{MA-PR.207}
\]
The comparison \(\log(3+|t|)\ll\ell_n\) follows
from the original selected-height intervals.
The count over \(\mathcal Z_n\) is bounded by
the full count (MA-PR.205), so no zero outside
the original window was silently discarded
from an upper bound.  Taking the two suprema
in (MA-PR.153) proves \(U_n=O(\log\nu_n)\).

\emph{3. Passage to the original whole-quotient norm.}
Equation (MA-PR.204) implies the quantitative
edge avoidance (MA-PR.151): for a fixed positive
constant \(c_4\),
\[
       \min(\beta,1-\beta)
                          \ge c_4(2+|\gamma|)^{-2}.
 \tag{MA-PR.208}
\]
There is no zero on either edge, and every
completed zero is inside the Hilbert strip.
Thus \(r_\rho=m_\rho\) for every \(\rho\),
for all fixed \(s\ge0\).
All factors are retained; nevertheless the
norm of \(\mathscr R_n\) on the whole \(Q_H\)
is still the norm that must be bounded.

The proof of MA-PR-012 uses the same original
outer contours
\(\Gamma_n=\partial([-1,2]+i[h_n,h_{n+1}])\),
but takes the actual Hilbert Mellin values
on the displaced inner vertical sides
\(\operatorname{Re}w=\delta_n,1-\delta_n\),
where \(\delta_n=\delta_0\nu_n^{-2}\).
By (MA-PR.208), \(\delta_0\) may be chosen
once so that these sides stay inside the
zero-free edge margins and the inner rectangle
contains exactly all zeros of the original
block.  For a general Hilbert vector, this
does not evaluate a Mellin transform on an
unjustified endpoint line.  The actual
Cauchy--Schwarz bound on those inner sides
is at most \(C\delta_n^{-1/2}\|h\|_H\);
the extra polynomial weight can only improve
this estimate.

The entire interpolant (MA-PR.162) consequently
has exactly the required full jets, and
(MA-PR.167) measures its class in the
unchanged \(J_n\) norm.  Equations
(MA-PR.154) and (MA-PR.207) yield
\[
      \|\iota_n\mathscr R_n\|_{Q_H\to J_n}
            \le C\nu_n^{D_0}e^{U_n}
            \le C'\nu_n^{D_1},
 \tag{MA-PR.209}
\]
which is (MA-PR.201).
Pole removal in \(\Xi\) has only specified
the existing nontrivial divisor.  It has
introduced neither a new Hilbert norm nor
a replacement block norm.

\emph{4. Actual returns and the trace identity.}
The no-proper-truncation condition of
MA-PR-010 now holds, and (MA-PR.209) is its
polynomial hypothesis in the exact original
norm.  Its actual construction
(MA-PR.108)--(MA-PR.109) proves the first
two identities, uniqueness, lift topology
and nuclear assertions in (MA-PR.202).
Since there are no discarded factors,
the correction \(E_g\) of (MA-PR.116) has
zero full jets at every completed zero.
The complete smooth jet-kernel theorem
therefore makes it the zero operator on
\(Q_\infty\).  This proves the last identity
in (MA-PR.202) and the actual common-series
trace equality (MA-PR.203).
\end{proof}

\paragraph{Limits of this arithmetic application.}
At centered widths \(\alpha>1/2\), the whole
critical strip already has a fixed interior
margin; that earlier case does not require
the endpoint argument.  At \(\alpha=1/2\)
there need not be a fixed margin, and
(MA-PR.204)--(MA-PR.207) supply the required
polynomial estimate.  At a fixed
\(0<\alpha<1/2\), the two Hilbert edges
are interior vertical lines.  The cited
outside-line count does not become
(MA-PR.205) there.

Indeed a disk of radius \(r\) centered at an
interior edge \(e+it\) is contained in a disk
centered at its nearest endpoint of radius
\(\operatorname{dist}(e,\{0,1\})+r\).
For sufficiently small fixed \(r\), the
resulting upper bound contains the term
\(C\operatorname{dist}(e,\{0,1\})\log\nu_n\).
Integration down to a generic
\(\nu_n^{-2}\) distance therefore retains
the \(O(\log^2\nu_n)\) loss of (MA-PR.157).
This identifies the limitation of this
particular argument; it asserts neither
arithmetic bad clusters nor failure of
polynomial costs at an interior edge.

The conclusion is about the actual smooth
and Hilbert quotients in (MA-PR.1).
It does not identify their maps with a
physical Fourier cutoff or extend an
original full Meyer common-space matrix.
In particular the narrow full-source
criteria of MA-MO-036 have a different
domain and remain separate.
The finite-length and nuclear-trace
placement is inherited from MA-PR-010--011
and BR-TR-01/02; no arithmetic positivity
or weight theorem is inferred from those
algebraic descriptions.

\paragraph{Primary references.}
Lagarias, J. C., Montgomery, H. L., \& Odlyzko, A. M.
(1979). A bound for the least prime ideal in the
Chebotarev density theorem.
\emph{Inventiones Mathematicae, 54}, 271--296.
\url{https://doi.org/10.1007/BF01390234}.
The invoked locators are Lemma~2.2,
(2.11)--(2.14), pp.~276--277, and Lemma~2.3,
(2.16)--(2.20), pp.~277--279.

Lagarias, J. C., \& Odlyzko, A. M. (1977).
Effective versions of the Chebotarev density theorem.
In A. Fr\"ohlich (Ed.), \emph{Algebraic number fields:
\(L\)-functions and Galois properties} (pp.~409--464).
Academic Press.  The imported analytic proof
interfaces are Section~5, Lemmas~5.1--5.6,
and the real-axis estimate following (8.3).

Thorner, J., \& Zaman, A. (2015).
\emph{Explicit results on the distribution of zeros
of Hecke \(L\)-functions} (arXiv:1510.08086).
\url{https://arxiv.org/abs/1510.08086}.
Lemma~2.5, printed p.~9, gives an explicit
version of the same radius-sensitive count.
